\subsection{THE LEVI-MALCEV THEOREM}

Let E be a complete normed vector space over R and u a continuous endomorphism of E. We have seen (Functions of a real variable, Chapter IV, § 2, no. 6) that the sequence \(\frac{u^{n}}{n!}\) is summable in \(\mathcal{L}(E)\) and we wrote

\[
e ^ {u} = \exp u = \sum_ {n = 0} ^ {\infty} \frac {u ^ {n}}{n !}.
\]

Now let E be a vector space over the field K and u a nilpotent endomorphism of E. The series \(\sum_{n=0}^{\infty}\frac{u^{n}}{n!}\) has only a finite number of non-zero terms and we can therefore write

\[
e ^ {u} = \exp u = \sum_ {n = 0} ^ {\infty} \frac {u ^ {n}}{n !}.
\]

This definition agrees with the above if K = R and if E is complete and normed. If v is another nilpotent endomorphism of E which commutes with u, then:

\[
\begin{array}{l} e ^ {u} e ^ {v} = \left(\sum_ {n = 0} ^ {\infty} \frac {u ^ {n}}{n !}\right) \left(\sum_ {p = 0} ^ {\infty} \frac {v ^ {p}}{p !}\right) = \sum_ {n, p = 0} ^ {\infty} \frac {u ^ {n} v ^ {p}}{n ! p !} \\ = \sum_ {q = 0} ^ {\infty} \frac {1}{q !} \left(\sum_ {n + p = q} \binom {q} {n} u ^ {n} v ^ {p}\right) = \sum_ {q = 0} ^ {\infty} \frac {1}{q !} (u + v) ^ {q} = e ^ {u + v}. \end{array}\tag{3}
\]

In particular, \(e^u e^{-u} = e^{-u}e^u = e^0 = 1\) and hence \(e^u\) is always an automorphism of E.

If further E is a (not necessarily associative) algebra and u is a (nilpotent) derivation of E, then \(e^{u}\) is an automorphism of the algebra E. For if \(x, y \in E\) then

\[
u ^ {p} (x y) = \sum_ {r + s = p} \binom {p} {r} u ^ {r} (x) u ^ {s} (y)
\]

for every integer \(p \geqslant 0\) (Leibniz's formula). It follows that

\[
\begin{array}{r l} e ^ {u} (x y) & = \sum_ {p \geqslant 0} \frac {1}{p !} u ^ {p} (x y) = \sum_ {p \geqslant 0} \sum_ {r + s = p} \frac {u ^ {r} (x)}{r !} \frac {u ^ {s} (y)}{s !} \\ & = \sum_ {r, s = 0} ^ {\infty} \frac {u ^ {r} (x)}{r !} \frac {u ^ {s} (y)}{s !} = e ^ {u} (x) e ^ {u} (y) \end{array}
\]

whence our assertion.

Now let g be a Lie algebra. If x belongs to the nilpotent radical of g, the derivation \(ad_{g}x\) of g is nilpotent. We can therefore make the following definition:

DEFINITION 6. A special automorphism of g is an automorphism of g of the form \(e^{ad_{\mathfrak{g}} x}\) , where x is in the nilpotent radical of g.

Clearly a special automorphism leaves every ideal of g stable.

DEFINITION 7. Let \(\mathfrak{g}\) be a Lie algebra and \(\mathfrak{r}\) its radical. A Levi subalgebra of \(\mathfrak{g}\) is any subalgebra of \(\mathfrak{g}\) supplementary to \(\mathfrak{r}\) .

A Levi subalgebra is isomorphic to \(\mathfrak{g} / \mathfrak{r}\) and hence is semi-simple. As a semi-simple subalgebra has only 0 in common with \(\mathfrak{r}\) , every semi-simple subalgebra \(\mathfrak{h}\) such that \(\mathfrak{g} = \mathfrak{r} + \mathfrak{h}\) is a Levi subalgebra; consequently the image of a Levi subalgebra under a surjective homomorphism is a Levi subalgebra.

THEOREM 5 (Levi--Malcev). A Lie algebra g always has a Levi subalgebra s. Every Levi subalgebra of g is the image of s under a special automorphism.

Let r denote the radical of g. We first treat two special cases.

\begin{enumerate}
\def\labelenumi{(\alph{enumi})}
\tightlist
\item
  \([g, r] = \{0\}\) .
\end{enumerate}

By Proposition 5, g is then the product of its centre r by Dg which is semi-simple. Hence Dg is a Levi subalgebra. Moreover, if \(s'\) is a semi-simple subalgebra, then \(s' = Ds'\) (Theorem 1), hence \(s' \subset Dg\) and Dg is the unique Levi subalgebra of g.

\begin{enumerate}
\def\labelenumi{(\alph{enumi})}
\setcounter{enumi}{1}
\tightlist
\item
  \([\mathfrak{g},\mathfrak{r}]\neq \{0\}\) and the only ideals of \(\mathfrak{g}\) contained in \(\mathfrak{r}\) are \(\{0\}\) and \(\mathfrak{r}\) .
\end{enumerate}

Then \([\mathfrak{g},\mathfrak{r}] = \mathfrak{r}\) , \([\mathfrak{r},\mathfrak{r}] = \{0\}\) and the centre of \(\mathfrak{g}\) is zero. Let M (resp. N) be the subspace of \(\mathcal{L}(\mathfrak{g})\) consisting of the linear mappings of \(\mathfrak{g}\) into \(\mathfrak{r}\) whose restriction to \(\mathfrak{r}\) is a homothety (resp. zero); N is therefore of codimension 1 in M. For \(m\in M\) , let \(\lambda (m)\) denote the ratio of the homothety of \(\mathfrak{r}\) defined by \(m\) . Let \(\sigma\) be the representation of \(\mathfrak{g}\) on \(\mathcal{L}(\mathfrak{g})\) canonically derived from the adjoint representation; recall that \(\sigma (x).u = [\mathrm{ad}_{\mathfrak{g}}x,u]\) for all \(x\in \mathfrak{g}\) and all \(u\in \mathcal{L}(\mathfrak{g})\) .

\[
\mathcal{L}(\mathfrak{g}) \supset \mathrm{M} \supset \mathrm{N} \supset \mathrm{P}
\]
Clearly \(\sigma (x)(\mathrm{M})\subset \mathrm{N}\) for all \(x\in \mathfrak{g}\) . Moreover, if \(x\in \mathfrak{r},y\in \mathfrak{g}\) and \(u\in \mathrm{M}\) , then\\
(4) \((\sigma (x).u)(y) = [x,u(y)] - u([x,y]) = -\lambda (u)[x,y]\) since \([\mathfrak{r},\mathfrak{r}] = \{0\}\) ; and (4) can be written:\\
(5) \(\sigma (x).u = -\mathrm{ad}(\lambda (u).x)\) .

As the centre of g is zero, the mapping \(x \mapsto ad_{g}x\) defines a bijection \(\phi\) of r onto a subspace P of \(\mathcal{L}(\mathfrak{g})\) . This subspace is stable under \(\sigma(\mathfrak{g})\) and contained in N since r is a commutative ideal and (5) shows that \(\sigma(x)(\mathrm{M}) \subset \mathrm{P}\) for \(x \in \mathfrak{r}\) . The representation of g on M/P = V derived from \(\sigma\) is therefore zero on r and defines a representation \(\sigma'\) of the semi-simple algebra g/r on V. For all \(y \in g/r\) , the space \(\sigma'(y)(V)\) is contained in N/P, which is of codimension 1 in V. Consequently (no. 2, Lemma 3) there exists \(u_{0} \in M\) such that \(\lambda(u_{0}) = -1\) and \(\sigma(x).u_{0} \in P\) for all \(x \in g\) . The mapping \(x \mapsto \phi^{-1}(\sigma(x).u_{0})\) is a linear mapping of g into r. By (5) its restriction to r is the identity mapping of r. Hence its kernel is a subspace s of g supplementary to r in g. As s is the set of \(x \in g\) such that \(\sigma(x).u_{0} = 0\) , s is a subalgebra of g and therefore a Levi subalgebra of g.

Let \(\mathfrak{s}'\) be another Levi subalgebra. For all \(x \in \mathfrak{s}'\) , let \(h(x)\) be the unique element of \(\mathfrak{r}\) such that \(x + h(x) \in \mathfrak{s}\) . Since \(\mathfrak{s}\) is a subalgebra and \(\mathfrak{r}\) is commutative, for \(x, y\) in \(\mathfrak{s}'\) :

\[
[ x + h (x), y + h (y) ] = [ x, y ] + [ x, h (y) ] + [ h (x), y ] \in \mathfrak {s}
\]

hence

\[
h ([ x, y ]) = (\operatorname{ad} x). h (y) - (\operatorname{ad} y). h (x).
\]

By Remark 2 of no. 2, there exists \(a \in \mathfrak{r}\) such that \(h(x) = -[x, a]\) for all \(x \in \mathfrak{s}'\) . Then:

\[
x + h (x) = x + [ a, x ] = (1 + \operatorname{ad} a). x.\tag{6}
\]

As \(\mathfrak{r}\) is commutative, \((\operatorname{ad} a)^2 = 0\) and hence \(1 + \operatorname{ad} a = e^{\operatorname{ad} a}\) . As \(\mathfrak{r} = [\mathfrak{s}, \mathfrak{r}]\) , \(e^{\operatorname{ad} a}\) is a special automorphism of \(\mathfrak{g}\) . By (6) this special automorphism maps \(\mathfrak{s}'\) to \(\mathfrak{s}\) .

\textbf{(c) General case:}

We argue by induction on the dimension \(n\) of the radical. There is nothing to prove if \(n = 0\) and hence it can be assumed that the theorem holds for Lie algebras whose radical is of dimension \(<\dim \mathfrak{r}\) . By (a) it suffices to consider the case where \([\mathfrak{g}, \mathfrak{r}] \neq \{0\}\) . As \([\mathfrak{g}, \mathfrak{r}]\) is nilpotent (no. 4, Proposition 6), its centre \(\mathfrak{c}\) is \(\neq \{0\}\) . Let \(\mathfrak{m}\) be a minimal non-zero ideal of \(\mathfrak{g}\) contained in \(\mathfrak{c}\) . If \(\mathfrak{m} = \mathfrak{r}\) , we have case (b). Therefore let \(\mathfrak{m} \neq \mathfrak{r}\) and let \(f\) be the canonical mapping of \(\mathfrak{g}\) onto \(\mathfrak{g}' = \mathfrak{g}/\mathfrak{m}\) . The radical of \(\mathfrak{g}'\) is \(\mathfrak{r}' = \mathfrak{r}/\mathfrak{m}\) . By the induction hypothesis, \(\mathfrak{g}'\) has a Levi subalgebra \(\mathfrak{h}'\) . Then \(\mathfrak{h} = f^{-1}(\mathfrak{h}')\) is a subalgebra of \(\mathfrak{g}\) containing \(\mathfrak{m}\) such that \(\mathfrak{h}/\mathfrak{m} = \mathfrak{h}'\) is semi-simple and hence having \(\mathfrak{m}\) as radical. By the induction hypothesis \(\mathfrak{h} = \mathfrak{m} + \mathfrak{s}\) where \(\mathfrak{s}\) is a semi-simple subalgebra. Then the equality \(\mathfrak{g}' = \mathfrak{r}' + \mathfrak{h}'\) implies \(\mathfrak{g} = \mathfrak{r} + \mathfrak{h} = \mathfrak{r} + \mathfrak{m} + \mathfrak{s} = \mathfrak{r} + \mathfrak{s}\) and hence \(\mathfrak{s}\) is a Levi subalgebra of \(\mathfrak{g}\).

Let \(\mathfrak{s}'\) be another Levi subalgebra of \(\mathfrak{g}\) . Then \(f(\mathfrak{s})\) and \(f(\mathfrak{s}')\) are two Levi subalgebras of \(\mathfrak{g}'\) and there exists by the induction hypothesis \(a' \in [\mathfrak{g}', \mathfrak{r}']\) such that \(e^{\mathrm{ad} a'}(f(\mathfrak{s}')) = f(\mathfrak{s})\) . If \(a \in [\mathfrak{g}, \mathfrak{r}]\) is such that \(f(a) = a'\) , it follows that:

\[
\mathfrak {s} _ {1} = e ^ {\mathrm{ad} a} \left(\mathfrak {s} ^ {\prime}\right) \subset \mathfrak {m} + \mathfrak {s} = \mathfrak {h}.
\]

Then \(\mathfrak{s}_1\) and \(\mathfrak{s}\) are two Levi subalgebras of \(\mathfrak{h}\) and by the induction hypothesis there exists \(b\in \mathfrak{m}\) such that \(e^{\mathrm{ad}b}(\mathfrak{s}_1) = \mathfrak{s}\) . Hence \(\mathfrak{s} = e^{\mathrm{ad}b}.e^{\mathrm{ad}a}(\mathfrak{s}')\) . Finally, as \(\mathfrak{m}\) is in the centre of \([\mathfrak{g},\mathfrak{r}]\) , \(e^{\mathrm{ad}b}.e^{\mathrm{ad}a} = e^{\mathrm{ad}(b + a)}\) and \(b + a\in [\mathfrak{g},\mathfrak{r}]\) , which completes the proof.

COROLLARY 1. Let \(\mathfrak{s}\) be a Levi subalgebra of \(\mathfrak{g}\) and \(\mathfrak{h}\) a semi-simple subalgebra of \(\mathfrak{g}\) .

\begin{enumerate}
\def\labelenumi{(\alph{enumi})}
\item
  There exists a special automorphism of \(\mathfrak{g}\) mapping \(\mathfrak{h}\) onto a subalgebra of \(\mathfrak{s}\) .
\item
  \(\mathfrak{h}\) is contained in a Levi subalgebra of \(\mathfrak{g}\) .
\end{enumerate}

Let \(\mathfrak{r}\) be the radical of \(\mathfrak{g}\) and \(\mathfrak{a} = \mathfrak{h} + \mathfrak{r}\) , which is a subalgebra of \(\mathfrak{g}\) . Then \(\mathfrak{a}/\mathfrak{r}\) is semi-simple and \(\mathfrak{r}\) is solvable, hence \(\mathfrak{r}\) is the radical of \(\mathfrak{a}\) and \(\mathfrak{h}\) is a Levi subalgebra of \(\mathfrak{a}\) . On the other hand, \(\mathfrak{a} \cap \mathfrak{s} = \mathfrak{h}'\) is a supplementary subalgebra to \(\mathfrak{r}\) in \(\mathfrak{a}\) and hence also a Levi subalgebra of \(\mathfrak{a}\) . Then there exists (Theorem 5) \(a \in [\mathfrak{a}, \mathfrak{r}]\) such that \(e^{\mathrm{ad} a}\) maps \(\mathfrak{h}\) onto \(\mathfrak{h}'\) . Now \(a \in [\mathfrak{g}, \mathfrak{r}]\) ; \(e^{\mathrm{ad} a}\) maps \(\mathfrak{h}\) onto a subalgebra of \(\mathfrak{s}\) and \(e^{-\mathrm{ad} a}(\mathfrak{s})\) is a Levi subalgebra of \(\mathfrak{g}\) containing \(\mathfrak{h}\) .

COROLLARY 2. For a subalgebra \(\mathfrak{h}\) of \(\mathfrak{g}\) to be a Levi subalgebra of \(\mathfrak{g}\) , it is necessary and sufficient that \(\mathfrak{h}\) be a maximal semi-simple subalgebra of \(\mathfrak{g}\) .

This follows immediately from Corollary 1.

COROLLARY 3. Let g be a Lie algebra and m an ideal of g such that g/m is semi-simple. Then g contains a subalgebra supplementary to m in g. In other words, every extension of a semi-simple Lie algebra is inessential.

Let \(\mathfrak{s}\) be a Levi subalgebra of \(\mathfrak{g}\) (Theorem 5). Its canonical image in \(\mathfrak{g}/\mathfrak{m}\) is a Levi subalgebra and therefore equal to \(\mathfrak{g}/\mathfrak{m}\) , hence \(\mathfrak{g} = \mathfrak{s} + \mathfrak{m}\) . Then an ideal of \(\mathfrak{s}\) supplementary in \(\mathfrak{s}\) to the ideal \(\mathfrak{m} \cap \mathfrak{s}\) is a subalgebra of \(\mathfrak{g}\) supplementary to \(\mathfrak{m}\) in \(\mathfrak{g}\) .

COROLLARY 4. Let g be a Lie algebra, r its radical, s a Levi subalgebra of g and m an ideal of g. Then m is the direct sum of m ∩ r which is its radical and m ∩ s which is a Levi subalgebra of m.

We know that \(\mathfrak{m} \cap \mathfrak{s}\) is the radical of \(\mathfrak{m}\) (§ 5, no. 5, Corollary 3 to Proposition 5). Let \(\mathfrak{h}\) be a Levi subalgebra of \(\mathfrak{m}\) and \(\mathfrak{s}'\) a Levi subalgebra of \(\mathfrak{g}\) containing \(\mathfrak{h}\) (Corollary 1). The algebra \(\mathfrak{m} \cap \mathfrak{s}'\) is an ideal of \(\mathfrak{s}'\) , is therefore semi-simple, and contains \(\mathfrak{h}\) and is therefore equal to \(\mathfrak{h}\) . Hence \(\mathfrak{m}\) is the direct sum of \(\mathfrak{m} \cap \mathfrak{r}\) and \(\mathfrak{m} \cap \mathfrak{s}'\) . There exists a special automorphism mapping \(\mathfrak{s}'\) onto \(\mathfrak{s}\) ; this automorphism leaves \(\mathfrak{r}\) and \(\mathfrak{m}\) invariant; hence \(\mathfrak{m}\) is the direct sum of \(\mathfrak{m} \cap \mathfrak{r}\) and \(\mathfrak{m} \cap \mathfrak{s}\) and \(\mathfrak{m} \cap \mathfrak{s}\) is a Levi subalgebra of \(\mathfrak{m}\) .

